\chapter{مفاهيم أساسية}
 
\subsection*{المعادلة التفاضلية \cite{diff_eqs_sols_apps}}
\addcontentsline{toc}{section}{المعادلة التفاضلية}
    هي علاقة بين المتغير المعتمد والمتغير المستقل تدخل فيها المشتقات أو التفاضلات وتسمى المعادلة التفاضلية اعتيادية إذا كان المتغير المعتمد دالة في متغير مستقل واحد وبالتالي لا تحتوي إلا على مشتقات عادية. حيث الشكل العام لها
\[
F(x,y',y'',y''',\dots,y^{(n)})=0
\]
حيث $y^{(n)}=\mfrac{d^ny}{dx^n}$ المشتقة من الرتبة $n$ للمتغير $y$ بالنسبة إلى $x$.

\begin{nexample}
    ليكن $x$ هو المتغير المستقل و $y$ هو المتغير المعتمد. فالعلاقات التالية تمثل معادلات تفاضلية اعتيادية
    \begin{align}
        &\frac{dy}{dx}+y=3\sin x\label{diff_eq_ex1}\\
        &\frac{d^2y}{dx^2}+\frac{1}{x}\frac{dy}{dx}+y=0\label{diff_eq_ex2}
    \end{align}
\end{nexample}

\subsection*{رتبة المعادلة التفاضلية \cite{diff_eqs_sols_apps}}
\addcontentsline{toc}{section}{رتبة المعادلة التفاضلية}
    إذا كانت المشتقة النونية $y^{(n)}$ هي أعلى مشتقة تظهر في المعادلة التفاضلية الاعتيادية قيل أن هذه المعادلة من الرتبة $n$.

%\begin{example}
%    المعادلة \eqref{diff_eq_ex1} هي معادلة اعتيادية من الرتبة الأولى. أما المعادلة \eqref{diff_eq_ex2} هي معادلة تفاضلية اعتيادية من الرتبة الثانية.
%\end{example}

\subsection*{درجة المعادلة التفاضلية \cite{diff_eqs_sols_apps}}
\addcontentsline{toc}{section}{درجة المعادلة التفاضلية}
    هي الاس المرفوع إليها أعلى مشتقة تظهر في المعادلة التفاضلية ، وقبل تحديد درجة المعادلة التفاضلية يجب وضعها في أبسط صورة قياسية صحيحة من حيث المشتقات.
\newpage
\begin{nexample}
	اعتبر المعادلة التفاضلية البسيطة:
	\[
	\frac{dy}{dx} + y = 0
	\]
	هذه معادلة تفاضلية من الرتبة الأولى لأن أعلى مشتقة فيها هي $\frac{dy}{dx}$، 
	وهي من الدرجة الأولى لأن أسّ أعلى مشتقة هو $1$.\\
	\noindent
\textbf{	مثال آخر:}
	\[
	\left(\frac{d^2 y}{dx^2}\right)^2 + \frac{dy}{dx} + y = 0
	\]
	هذه معادلة من الرتبة الثانية لأن أعلى مشتقة هي $\frac{d^2 y}{dx^2}$،
	ومن الدرجة الثانية لأن هذه المشتقة مرفوعة للأس $2$.
\end{nexample}

\subsection*{المعادلة التفاضلية الخطية \cite{diff_eqs_sols_apps}}
\addcontentsline{toc}{section}{المعادلة التفاضلية الخطية}
    هي المعادلة الخطية في المتغير المعتمد و مشتقاته جميعاً. الصورة العامة للمعادلة التفاضلية الخطية من الرتبة $n$ 
    \[
    p_n(x)y^{(n)}+p_{n-1}(x)y^{(n-1)}+\dots+p_1(x)y'+p_0(x)y=Q(x)
    \]
    أو
    \[
    \sum_{i=0}^{n}p_i(x)y^{(i)}=Q(x)
    \]
    حيث المتغير المعتمد $y$ و جميع مشتقاته مرفوعة إلى الاس واحد ولا توجد حواصل ضرب مشتركة في ما بينها ، و الدوال $Q(x)$ و $p_i(x)$ هي دوال للمتغير المستقل $x$ خطية أم غير خطية لا تؤثر على خطية المعادلة التفاضلية.
\vskip10pt
\begin{note}
    من تعريف المعادلة التفاضلية الخطية يمكننا القول أي معادلة تفاضلية فيها المتغير المعتمد $y$ أو أحد مشتقاته مرفوعة إلى اس غير الواحد أو وجدنا حاصل ضرب في ما بينها. سوف نطلق عليها معادلة تفاضلية غير خطية.
\end{note}
\newpage
\begin{nexample}
    لدينا المعادلة
    \[
    x^2y^{''}+xy'+y=\cos x
    \]
    معادلة تفاضلية اعتيادية خطية. أما المعادلة
    \[
    yy^{'}+y^{''}=e^{3x}
    \]
    هذه المعادلة التفاضلية غير خطية لوجود حاصل الضرب $yy^{'}$. وأيضا لدنيا
\[
y^{'}+x\sqrt[4]{y}=\cot x
\]
وهذه كذلك ليست خطية فيها المتغير المعتمد $y$ مرفوع إلى الاس $1/4$ (غير الواحد).
\end{nexample}

\subsection*{المعادلة التفاضلية الخطية المتجانسة \cite{diff_eqs_sols_apps}}
\addcontentsline{toc}{section}{المعادلة التفاضلية الخطية المتجانسة}
اذا انعدمت الدالة $Q(x)$ (اي ان $Q(x) = 0$) من المعادلة التفاضلية
\[
p_n(x)y^{(n)}+p_{n-1}(x)y^{(n-1)}+\dots+p_1(x)y'+p_0(x)y=Q(x)
\]
أو
\[
\sum_{i=0}^{n}p_i(x)y^{(i)}=Q(x)
\]
لجميع قيم $x$ فأنها تصبح معادلة خطية متجانسة ، عدا ذلك تكون المعادلة غير متجانسة

\begin{nexample}
	المعادلة
	\[
	x^2 y'' + xy' + (x^2 - 1)y = 0
	\]
	هي معادلة تفاضلية خطية متجانسة من الرتبة الثانية. بينما المعادلة
	\[
	xy' + (\sin x)y = x^2(\sin x + 2)
	\] 
	هي معادلة تفاضلية خطية غير متجانسة من الرتبة الاولى.
\end{nexample}

\newpage

\subsection*{حل المعادلة التفاضلية \cite{diff_eqs_pt1}}
\addcontentsline{toc}{section}{حل المعادلة التفاضلية}
تسمى الدالة $y=y(x)$ حلاً للمعادلة التفاضلية $F(x,y,y',y'',\dots,y^{(n)})$ إذا كانت
\begin{enumerate}
    \item قابلة للاشتقاق $n$ من المرات
    \item تحقق المعادلة التفاضلية أي أن $F(x,y(x),y'(x),y''(x),\dots,y^{(n)}(x))$
\end{enumerate}


\begin{nexample}
    أثبت أن $y(x)=c \sin x$ حلاً للمعادلة التفاضلية $y''+y=0$ حيث $c$ ثابت.
\end{nexample}
\begin{solution}
    نشتق العلاقة ونقوم بتعويضها في المعادلة التفاضلية
\[
y=c\sin x\Rightarrow y'=c\cos x\Rightarrow y''=-c\sin x
\]
الآن نعوض
\[
y''+y=-c\sin x+c\sin x=0
\]
\end{solution}

\section*{انواع حلول المعادلات التفاضلية \cite{diff_eqs_pt1}}
\addcontentsline{toc}{section}{انواع حلول المعادلة التفاضلية}
\subsection*{الحل العام}
 للمعادلة التفاضلية من الرتبة $n$ هو حل يحتوي على $n$ من الثوابت الاختيارية وبالطبع يحقق المعادلة التفاضلية.
\subsection*{الحل الخاص}
 هو أي حل يحقق المعادلة التفاضلية لا يشتمل على أي ثوابت اختيارية وقد نحصل عليه احيانا بالتعويض عن الثوابت الاختيارية في الحل العام بقيم محددة.
\newpage
\begin{nexample}
    الحل العام للمعادلة التفاضلية $y'''-5y''+6y'=0$ يمكن ايجادمن خلال المعادلة المميزة للمعادلة التفاضلية الناتجة من تحويل شكل المعادلة الى صيغة المؤثر
    \[
    (D^3 - 5D^2 + 6D) y = 0
    \]
    اذن المعادلة المميزة تكون
    \[
    m^3 - 5m^2 + 6m = 0
    \]
    نلاحظ يمكننا ان نحلل الطرف الايسر بالعامل المشترك
    \[
    m(m^2 - 5m + 6) = 0
    \]
    الآن بالتجربة
    \[
    m(m-2)(m-3) = 0
    \]
    اذن الحلول للمعادلة المميزة
    \[
    m_1 = 0 , m_2 = 2, m_3 = 3
    \]
    بالتالي الحل العام يكون
    \[
    y = C_1 e^{m_1x} + C_2 e^{m_2 x} + C_3 e^{m_3 x}
    \]
   حيث $C_1, C_2, C_3 $ ثوابت اختيارية ، اذن الحل العام النهائي
    \[
    y(x) = C_1 + C_2 e^{2x} + C_3 e^{3x}
    \]
     ويمكننا ان نختار قيم محددة للثوابت الاختيارية ويتحول الحل العام الى حل خاص ، مثلاً نختار 
     $C_1 = 1, C_2, -4, C_3 = 3$ ، ليصبح الحل
    \[
    y(x) = 1 -4e^{2x} + 3e^{3x}
    \]
\end{nexample}
